%% =============================================================================
%%  config.tex — caderno de estudo de livro
%% =============================================================================

% ---- identificação do caderno ------------------------------------------------
\newcommand\ebTitulo{«TITULO»}
\newcommand\ebSubtitulo{Caderno de estudo do livro, lição a lição}
\newcommand\ebIncipit{ἀνακρίνοντες τὰς γραφάς}   % At 17:11 “examinando as Escrituras”
\newcommand\ebSerie{«SERIE»}
\newcommand\ebNumero{«NUMERO»}
\newcommand\ebAutor{«AUTOR»}
\newcommand\ebData{«DATA»}
\newcommand\ebEpigrafe{«EPIGRAFE»}
\newcommand\ebEpigrafeRef{«EPIGRAFE_REF»}

% ---- o livro estudado ---------------------------------------------------------
\newcommand\livroNome{«TITULO»}
\newcommand\livroEditora{}
\newcommand\livroEdicao{}
\newcommand\livroDados{dados/livros/«SIGLA»}   % licoes.tsv e textos.tsv
\newcommand\livroUnidade{Lição}            % Lição | Capítulo | Artigo ...
\newcommand\livroUnidadePlural{Lições}
\newcommand\livroParte{Parte}
% impresso: linhas e caixinhas para escrever à mão | digital: só licoes/NN.tex
\newcommand\livroModo{impresso}
\newcommand\livroLinhas{5}
% Lições a incluir: vazio = todas; ex.: 1-12 ou 1,4,7
\newcommand\livroLicoes{}

% ---- aparência ----------------------------------------------------------------
\newcommand\ebTema{«TEMA»}          % lapis | purpura | oliveira | sepia
\newcommand\ebPapel{«PAPEL»}         % a4 | b5
\newcommand\ebFundoPagina{nao}
\newcommand\ebAbas{sim}
\newcommand\ebAbasPor{parte}
\newcommand\ebInterlinearModo{estudo}
\newcommand\ebAcentosHebraicos{nao}
\newcommand\ebFontesPersonalizadas{}
